\documentclass[12pt,a4paper,fontset=none]{ctexart}
\usepackage{geometry}
\geometry{top=2cm,bottom=2cm,left=2.4cm,right=2.4cm}
\usepackage{xcolor}
\definecolor{ink}{HTML}{1F2D3D}
\definecolor{accent}{HTML}{8C2A2A}
\definecolor{rule}{HTML}{C9A24B}
\definecolor{gray}{HTML}{9AA7B4}

% 候选字体
\newCJKfontfamily\heiR{Noto Sans CJK SC}            % 思源黑体 正
\newCJKfontfamily\heiB{Noto Sans CJK SC}[BoldFont={Noto Sans CJK SC Bold}]
\newCJKfontfamily\songR{Noto Serif CJK SC}          % 思源宋体 正
\newCJKfontfamily\songB{Noto Serif CJK SC}[BoldFont={Noto Serif CJK SC Bold}]
\newCJKfontfamily\kaiR{LXGW WenKai}                 % 霞鹜文楷 正
\newCJKfontfamily\kaiM{LXGW WenKai Medium}          % 霞鹜文楷 中

\setCJKmainfont{Noto Sans CJK SC}
\linespread{1.5}
\setlength{\parindent}{2em}

\newcommand{\sep}{\par\vspace{4pt}{\color{rule}\rule{\linewidth}{0.6pt}}\par\vspace{10pt}}

\begin{document}
\thispagestyle{empty}

\begin{center}
{\heiB\bfseries\Large\color{ink} 中文字体对比样张}\\[2pt]
{\heiR\small\color{gray} 同一段文字 · 三种字体 · 供挑选}
\end{center}
\vspace{6pt}
{\color{rule}\rule{\linewidth}{1pt}}
\vspace{14pt}

% ---------- 1 思源黑体 ----------
{\heiR
{\heiB\bfseries\large\color{accent} 一、思源黑体（Noto Sans CJK SC）—— 当前所用 · 苹方风格}\par
\vspace{4pt}
{\small\color{gray} 现代无衬线 · 屏幕与封面友好 · 干净利落，偏“产品/科技”气质}\par
\vspace{6pt}
{\heiB\bfseries\large 第六章\quad 康德的哥白尼革命}\par
休谟的炸弹惊醒了一个德国人康德。他憋出了人类思想史上最艰深也最伟大的一招：不是我们的认识符合对象，而是\textbf{对象符合我们的认识}。我们永远看不到“物自体”，只能看到经过眼镜加工的“现象”。\quad{\color{accent}“我思故我在。”}
}
\sep

% ---------- 2 思源宋体 ----------
{\songR
{\songB\bfseries\large\color{accent} 二、思源宋体（Noto Serif CJK SC）—— 经典书卷气}\par
\vspace{4pt}
{\heiR\small\color{gray} 衬线/宋体 · 长文阅读耐看 · 印刷书籍质感，偏“典籍/学术”气质}\par
\vspace{6pt}
{\songB\bfseries\large 第六章\quad 康德的哥白尼革命}\par
休谟的炸弹惊醒了一个德国人康德。他憋出了人类思想史上最艰深也最伟大的一招：不是我们的认识符合对象，而是\textbf{对象符合我们的认识}。我们永远看不到“物自体”，只能看到经过眼镜加工的“现象”。\quad{\color{accent}“我思故我在。”}
}
\sep

% ---------- 3 霞鹜文楷 ----------
{\kaiR
{\kaiM\large\color{accent} 三、霞鹜文楷（LXGW WenKai）—— 温润文学感（推荐用于故事）}\par
\vspace{4pt}
{\heiR\small\color{gray} 楷体风 · 手写温度 · 最适合散文/故事正文，偏“人文/书信”气质}\par
\vspace{6pt}
{\kaiM\large 第六章\quad 康德的哥白尼革命}\par
休谟的炸弹惊醒了一个德国人康德。他憋出了人类思想史上最艰深也最伟大的一招：不是我们的认识符合对象，而是{\kaiM 对象符合我们的认识}。我们永远看不到“物自体”，只能看到经过眼镜加工的“现象”。\quad{\color{accent}“我思故我在。”}
}
\sep

% ---------- 小字测试（验证可读性）----------
{\heiB\bfseries\color{ink} 小字号可读性（10pt，同句三体对照）：}\par
\vspace{4pt}
{\small\heiR 黑体：道可道，非常道；名可名，非常名。无名天地之始，有名万物之母。}\par
{\small\songR 宋体：道可道，非常道；名可名，非常名。无名天地之始，有名万物之母。}\par
{\small\kaiR 文楷：道可道，非常道；名可名，非常名。无名天地之始，有名万物之母。}\par

\vfill
{\color{rule}\rule{\linewidth}{1pt}}\par\vspace{4pt}
{\heiR\small\color{gray} 建议：故事/散文用「霞鹜文楷」或「思源宋体」更有书卷气；旗舰范本（含表格/术语/代码感）保持「思源黑体」更清晰。可两份文档用不同字体。}
\end{document}
